% ================================
% === Reviewer 2 ===
% ================================
\reviewerblock{Reviewer~2}{%
We thank the reviewer for acknowledging that the manuscript is well-written, technically sound, and empirically validated, and for raising the central novelty concern.}

\cresp{The major drawback is related to insufficient novelty of the technical contribution. Most of the contributions proposed in the paper relate to previously discussed topics within vision-language modeling, semantic alignment, and uncertainty-aware learning. Although the idea of distinguishing between semantic uncertainty and supervision uncertainty is quite appealing, it does not constitute a new theoretical principle or algorithmic idea. It appears that the results are based mainly on the correct combination of existing methods rather than the development of a novel approach itself.}
{We acknowledge that the original presentation did not distinguish the methodological contribution clearly enough from a combination of existing components. We therefore change the framing, formulation, and empirical validation together. In the Introduction, we rewrite the gap and contribution paragraphs so that the contribution is the functional separation of two measurable failure modes, rather than the mere use of a semantic prior and distillation. In the Method, we add a new subsection, ``Functionally Decoupled Ambiguity Modeling.'' Its diagnostic paragraph reports that the baseline Top-10 set already covers \res{$83.9\%$} of positive labels and retains \res{$67.8\%$} of its false negatives as locally recoverable candidates, indicating that much of the remaining semantic error is local ranking and calibration rather than complete retrieval failure. We then add Eq.~\eqref{eq:risk_decomp}, which separates supervisory ambiguity as a label-side conditional-entropy term from semantic ambiguity as a feature-side inter-class aliasing term. The text immediately following the equation maps these two terms to UTD's supervision stabilization and SLR-C's local ranking/calibration, making the decoupling principle explicit and testable rather than post hoc.

We test this principle in two newly added ablation blocks. ``Unified versus Functionally Decoupled Ambiguity Handling'' preserves the component families but changes their organization: under identical features and selection rules, merging the semantic prior and rationale teacher into one target is weaker than decoupled FDIL (\res{$55.66$ vs.\ $57.49$ Avg F1}), although the strongest single-signal control remains close. ``Negative Controls on the Text Teacher'' replaces matched rationales with shuffled or generic text and modifies the agreement gate; shuffling rationales reduces Macro-F1 by \res{$7.8$} and Hard-F1 by \res{$7.6$}, while removing or flattening the gate has smaller effects. Finally, we revise the Discussion to state the bounded conclusion: decoupling gives a modest aggregate benefit with separately measurable and deployable functions, but it is not the only competitive organization.
\revloc{the rewritten gap and contribution paragraphs in Sec.~\ref{sec:intro}; the new Functionally Decoupled Ambiguity Modeling subsection, diagnostic paragraph, and Eq.~\eqref{eq:risk_decomp} in Sec.~\ref{sec:decoupling_principle}; the new Unified versus Functionally Decoupled Ambiguity Handling and Negative Controls subsections with Tables~\ref{tab:unified_decoupled} and~\ref{tab:negative_controls}; the bounded interpretation in Sec.~\ref{sec:discussion}}
}
\showUnifiedDecoupled
\showNegativeControls
